\section{Fiscal: what does it cost, and is it paid for?}
\label{sec:fiscal}

The fourth question decides whether the bill can move at all. A member who cannot
explain the score will not carry the bill, and a bill that trips a budget point of
order does not reach the floor. H. R. 9510 is written to be scored cleanly.

The bill only authorizes; it does not appropriate, and it creates no entitlement. It
authorizes appropriations to the FDA Salaries and Expenses account of 18 million
dollars in fiscal year 2027, 12 million in 2028, 10 million in 2029, and 9 million in
each of 2030 and 2031, for a five-year total of 58 million dollars, with estimated
outlays of 53.5 million dollars inside the window and the small remainder after it
\cite{Kawchak2026HR9510Bill}. Crucially for scorekeeping, the bill creates no direct
spending and no revenue effect, so it carries no mandatory cost under the Statutory
Pay-As-You-Go Act and adds nothing to the House CutGo scorecard
\cite{paygo2010, houserules2025}. The only fee change is a carve-out: a verification
record filed on its own is not a submission and bears no user fee, and the bill
leaves the pricing of any incremental review workload to the next medical-device
user-fee reauthorization rather than imposing a new statutory fee. \autoref{tab:fiscal}
sets out the schedule, and a formal Congressional Budget Office estimate would follow
introduction in the ordinary course \cite{cbo2024costestimates}.

\needspace{12\baselineskip}\begin{table}[H]
\footnotesize
\begin{tabularx}{\textwidth}{@{}L{1.6cm} C{6cm} Y@{}}
\toprule
\textbf{Fiscal year} & \textbf{Authorization} & \textbf{Principal activity}\\
\midrule
FY 2027 & \$18 million & Rulemaking under section 515D(h)\\
FY 2028 & \$12 million & Verification-records system build\\
FY 2029 & \$10 million & Steady-state operation and inspection\\
FY 2030 & \$9 million & Steady-state operation\\
FY 2031 & \$9 million & Steady-state, then transition to user fees\\
\midrule
Total & \$58 million & \$53.5 million outlaid within the window; no direct spending\\
\bottomrule
\end{tabularx}
\caption{The authorization schedule of section 5 and its budgetary effect: a bounded,
declining authorization with no net mandatory spending.}
\label{tab:fiscal}
\end{table}

The decisive number for a skeptical member is the cost asymmetry made concrete. A
single cleared interaction class costs about 4,420 dollars to verify, of which the
ten-gate compute is roughly 3,710 dollars, while the conventional path of manual code
review and serial validation runs near 3 million dollars in a program-year; the
automated gate operates the same program-year for about 155,000 dollars, a reduction
of roughly nineteen to one \cite{Kawchak2026HR9510Bill,
Kawchak2025AcceleratingFDAComplianceCost}. \autoref{fig:fiscal} places the gate's
small, predictable cost against the unbounded cost of the failure it prevents.

\begin{psgfig}
\node[psgevid] (g) at (0,0) {Gate suite\\\$3,710 per run};
\node[psgtwo] (c) at (3.9,0) {Per interaction\\\$4,420};
\node[psgone] (a) at (7.8,0) {Automated year\\\$155,000};
\node[psgcaut] (f) at (7.8,-2.0) {Conventional year\\\$3,000,000};
\draw[psgedge] (g)--(c); \draw[psgedge] (c)--(a);
\draw[psgedge] (a)--(f) node[psglabel,midway,right]{19 to 1};
\end{psgfig}
\vspace{0.15cm}
\psgcaption{\textbf{Figure 5.} The cost asymmetry: a small, fixed verification cost
against the unbounded cost it removes, a roughly nineteen-to-one reduction over
conventional review.}
\label{fig:fiscal}

Fiscal is answered with a clean score: a bounded authorization, no mandatory
spending, no new fee, and a verification cost an order of magnitude below the status
quo, which is precisely what a member needs to explain to the Budget Committee.
